% proofs/09_regret.tex
% The benchmark-policy one-step deviation bound.
% Fragment, to be \input by appendix.tex.

\section{The benchmark policy: the interim regret bound}
\label{sec:regret}

At each calibration node, a one-step deviation changes the blockholder's plan at one signal while
the benchmark price system and beliefs stay fixed. The result below bounds the largest gain from
that deviation on the stated signal support.

\csname unlabelledtrue\endcsname
\begin{proposition}[Benchmark-policy interim regret]
\label{prop:regret}
Assume standing conditions (S1) to (S11). The fixed pass in (R-3) also satisfies the common noise
channel, control-node convention, and identified flagged types in (S12) to (S14). The following
additional conditions are in force.
\begin{enumerate}[label=(R-\arabic*),leftmargin=3.2em,itemsep=3pt]
\item \emph{Calibration.} The benchmark cutoffs are
  \[
  k^b=(0.9425017266871091, 1.8484512098302512),
  \]
  liquidity is $\kappa=0.5$, order size is two noise lumps, and $H=10$. The nodes pair
  $T\in\{3,5,10\}$ with
  \[
  \tau\in\left\{
  \begin{aligned}
  &0.09239820387429526, 0.09346755804663053, 0.09453534811956685,\\
  &0.09565657882778708, 0.09703177146201895
  \end{aligned}
  \right\},
  \]
  and the remaining primitives take their benchmark calibration values. In particular,
  $m_1>0$, $\bar b>b_0$, $\sigma_s>0$, $\sigma_\xi>0$, $\beta>0$, and $\chi>0$.
\item \emph{Signal support.} The signal has its calibrated Gaussian law restricted to
  $I_s=[s_{\mathrm{lo}},s_{\mathrm{hi}}]$, where
  $s_{\mathrm{lo}}=\mu_v-6\sigma_s$ and $s_{\mathrm{hi}}=\mu_v+6\sigma_s$. Its density is divided by the
  truncation constant
  \[
  Z_s=\Prb(s\in I_s)
  =\Phi\!\left(\frac{s_{\mathrm{hi}}-\mu_v}{\sigma_s}\right)
  -\Phi\!\left(\frac{s_{\mathrm{lo}}-\mu_v}{\sigma_s}\right)>0.
  \]
  The density is continuous, so every finite breakpoint set is null.
\item \emph{Fixed pooled pass and deviation convention.} At each node, the pooled pricing pass is
  computed once at the benchmark cutoffs with the run-up path present and is then held fixed. A
  deviation changes only the deviating signal's plan. Benchmark prices and beliefs do not react.
  Every off-path pooled history uses the fixed reference belief of the mark-path type that can
  produce it. At a cutoff the benchmark assigns Hold at the Exit-Hold cutoff and Voice at the
  Hold-Voice cutoff, so Exit is assigned for $s<k_1^b$, Hold for
  $k_1^b\le s<k_2^b$, and Voice for $s\ge k_2^b$.
\end{enumerate}

For calibration node $n$, let $U_{j,n}(s)$ be the payoff from plan
$j\in\{\mathsf{E},\mathsf{H},\mathsf{V}\}$ under (R-3), and let $j^b(s)$ be the benchmark plan. Define
\begin{equation}
\label{eq:regret-definition}
R_n(s)=\max_{j\in\{\mathsf{E},\mathsf{H},\mathsf{V}\}}U_{j,n}(s)-U_{j^b(s),n}(s).
\end{equation}
Use the benchmark payoff at the Hold-Voice cutoff as the payoff normaliser,
\begin{equation}
\label{eq:regret-normaliser}
N_n:=U_{\mathsf{V},n}(k_2^b)>0.
\end{equation}
The reported certificate value $\overline R_n$ satisfies
\begin{equation}
\label{eq:regret-bound}
\operatorname*{ess\,sup}_{s\in I_s}R_n(s)\le\overline R_n,
\qquad
\operatorname*{ess\,sup}_{s\in I_s}\frac{R_n(s)}{N_n}
\le\frac{\overline R_n}{N_n}.
\end{equation}
The record reports the bound at every node.
\end{proposition}
\csname unlabelledfalse\endcsname

\begin{proof}
Fix a node and suppress its subscript. Form a partition of $I_s$ from the two benchmark cutoffs,
every jump of the Voice accumulation length $n(s)$ in $I_s$, every signal in $I_s$ at which a
threshold-crossing date changes, and the two support endpoints. Remove exact duplicates and sort the
points. On each open
piece $I=(l,r)$, the assigned plan, the Voice type, the crossing date, the filing date, and the
filing status are constant. Each payoff formula on that piece extends continuously to its closure
from the relevant side. The partition is finite, and its endpoints are null by (R-2). It therefore
suffices to bound every branch closure.

Write
\[
b=b^*(s),
\qquad
v_h=\mu_v+\beta(s-\mu_v),
\qquad
C=C_0\exp\!\left[-\frac{\chi(s-\mu_v)}{\sigma_s}\right].
\]
For Voice type $m$, let $e_m$ be the fixed pass's expected control-node entry probability,
$q_m$ its expected product of entry and control-node price, and $P_{dm}$ its expected pooled
execution price at date $d$. Set
\[
A_m=\frac1m\sum_{d=0}^{m-1}P_{dm}.
\]
Direct substitution into the blockholder's payoff gives the three branch formulas
\begin{align}
\label{eq:regret-branches-pooled}
U_{\mathsf{E}}&=b_0P_{00},
\nonumber\\
U_{\mathsf{H}}&=b_0\bigl((1-e_0)v_h+q_0+e_0m_0\bigr),
\\
Y_m&=(1-e_m)(v_h+\Delta_V)+q_m+e_m m_1,
\nonumber\\
U_{\mathsf{V}}&=bY_m-A_m(b-b_0)-C
\quad\text{on a pooled Voice branch}.
\nonumber
\end{align}
The Exit payoff is constant and the Hold payoff is affine. On a flagged Voice branch with filing
date $f$, put
\[
r_f=\min\{1,(f+1)/m\},
\qquad
B^F=b_0+r_f(b-b_0),
\qquad
A_{mf}=\frac1m\sum_{d=0}^{\min\{f,m-1\}}P_{dm}.
\]
At engagement posterior one, the flagged price satisfies
\[
P_F=(1-p_F)(v_h+\Delta_V)+p_F(P_F+m_1).
\]
This fixed point makes expected terminal value equal to $P_F$. The flagged Voice payoff therefore
reduces to
\begin{equation}
\label{eq:regret-branch-flagged}
U_{\mathsf{V}}=B^FP_F-A_{mf}(b-b_0)-C.
\end{equation}
Thus every payoff is given by a closed form on each branch, apart from the unique scalar price root
in \eqref{eq:regret-branch-flagged}, whose uniqueness follows from Lemma~\ref{lem:pricing-root}.

It remains to cover the space between mesh points. The target and engagement cost satisfy
\[
b'(s)=\frac{\bar b-b_0}{2\sigma_s}
\left(1+\left(\frac{s-\mu_v}{\sigma_s}\right)^2\right)^{-3/2},
\qquad
C'(s)=-\frac{\chi}{\sigma_s}C(s).
\]
Let $b'_{\max,I}$ be the maximum of $b'$ on the closed piece, attained at the point nearest
$\mu_v$. Differentiating the pooled Voice branch gives
\[
U_{\mathsf{V}}'=b'(Y_m-A_m)+b(1-e_m)\beta+\frac{\chi}{\sigma_s}C.
\]
Since $Y_m-A_m$ is affine, a valid Lipschitz constant on $I$ is
\begin{equation}
\label{eq:regret-LP}
L_{\mathsf{V}}^{P}
=b'_{\max,I}\max_{s\in\{l,r\}}|Y_m(s)-A_m|
+\bar b(1-e_m)\beta+\frac{\chi}{\sigma_s}C(l).
\end{equation}
The other two constants are $L_{\mathsf{E}}=0$ and
$L_{\mathsf{H}}=b_0(1-e_0)\beta$.

For the flagged branch, put $V=v_h+\Delta_V$ and
\[
a(P)=1-p(P)=\Phi\!\left(\frac{P+K+m_1-\bar S}{\sigma_\xi}\right),
\qquad a=a(P_F).
\]
Implicit differentiation of the price equation gives
\begin{equation}
\label{eq:regret-price-derivative}
\frac{dP_F}{ds}
=\beta\frac{a^2}{a^2+(m_1/\sigma_\xi)\phi_{\!N}(u)},
\qquad
0<\frac{dP_F}{ds}\le\beta,
\end{equation}
where $u=(P_F+K+m_1-\bar S)/\sigma_\xi$ and $\phi_{\!N}$ is the standard normal density. The root
also satisfies
\[
V\le P_F\le V+m_1\frac{1-a(V)}{a(V)}.
\]
Since $V$ rises in $s$ while $(1-a(V))/a(V)$ falls, the endpoint mixture
\[
P_{\mathrm{abs},I}
=\max\left\{|V(l)|,
\left|V(r)+m_1\frac{1-a(V(l))}{a(V(l))}\right|\right\}
\]
bounds $|P_F|$ on $I$. Differentiating \eqref{eq:regret-branch-flagged} then gives the valid constant
\begin{equation}
\label{eq:regret-LF}
L_{\mathsf{V}}^{F}
=b'_{\max,I}\bigl(r_fP_{\mathrm{abs},I}+|A_{mf}|\bigr)
+\bar b\beta+\frac{\chi}{\sigma_s}C(l).
\end{equation}

On each piece take $L_{\mathsf{V}}=L_{\mathsf{V}}^P$ or
$L_{\mathsf{V}}=L_{\mathsf{V}}^F$ according to its filing status.
Let $a_I$ be the benchmark plan on a piece and let
$G_{j,a_I}=U_j-U_{a_I}$ for either alternative $j$. The preceding bounds make
$G_{j,a_I}$ Lipschitz on the branch closure with constant
$L_{j,a_I}=L_j+L_{a_I}$. Divide $[l,r]$ into
\[
n_I=\left\lceil\frac{r-l}{10^{-5}}\right\rceil
\]
equal segments, and call the endpoint mesh $\mathcal G_I$. Every point of the piece lies within
$\delta_I=(r-l)/(2n_I)$ of a mesh point. Hence
\begin{equation}
\label{eq:regret-cover}
\sup_{s\in I}G_{j,a_I}(s)
\le \max_{s\in\mathcal G_I}G_{j,a_I}(s)+L_{j,a_I}\delta_I.
\end{equation}
The certificate is rigorous in exact arithmetic. It adds $5\times10^{-12}$ to the right side for
floating-point evaluation and the reported price-root residual. Define $\overline R$ as the maximum
of zero and these adjusted bounds over every piece and
both alternative plans. Equation~\eqref{eq:regret-cover} then bounds $R(s)$ on every open piece.
The omitted endpoints form a null set, so it also bounds the essential supremum on $I_s$. This is
the first inequality in \eqref{eq:regret-bound}. Dividing by the positive normaliser $N$ proves the
second.
\end{proof}

The bound establishes nothing about equilibrium existence, order size, speed, or a reacting
blockholder.
